\section{Related Work}
\label{sec:related}

\paragraph{Volatility forecasting and strong baselines.}
Volatility, unlike returns, is persistently forecastable, and a mature econometric literature defines the bar any new predictor must clear. The GARCH family models conditional variance parametrically \citep{engle1982arch,bollerslev1986garch,nelson1991egarch}; a parallel tradition treats volatility as approximately observable through realised measures \citep{andersen2003modeling}, of which Corsi's heterogeneous autoregressive (HAR) model is the parsimonious workhorse \citep{corsi2009har}. These simple models are hard to beat out of sample \citep{hansen2005garch,poon2003forecasting}, and rich macro-financial predictor sets add little once a strong autoregressive benchmark is conditioned upon \citep{paye2012dejavol}. The credible question is therefore not whether a source predicts volatility in isolation, but whether it adds value \emph{beyond} a strong price baseline.

\paragraph{Text in finance.}
A large literature argues that financial text carries risk-relevant signal: media pessimism moves prices \citep{tetlock2007giving}, news sentiment predicts returns in recessions \citep{garcia2013sentiment}, finance-specific lexicons and learned term weights refine measurement \citep{loughran2011liability,loughran2016survey,jegadeesh2013word}, and market-level text proxies for risk perception and macroeconomic state \citep{manela2017nvix,bybee2024business}. Closest to us are the designs whose evaluation choices our protocol stress-tests: \citet{kogan2009predicting} regress firm volatility on 10-K word features; FinText forecasts realised volatility with financial word embeddings \citep{rahimikia2021fintext}; FETILDA fine-tunes long-document encoders on 10-K volatility regression \citep{fetilda2022}; and MDRM predicts post-earnings-call volatility from verbal and vocal cues \citep{qin2019mdrm}. None imposes a \emph{recalibrated} price reference, a firm-identity control, or day-clustered inference; baselines are often a single lagged feature rather than HAR-RV, panels short, and significance against the benchmark rarely tested.

\paragraph{Neural and long-document modelling.}
Transformer encoders are the strongest general text representations \citep{vaswani2017attention,devlin2019bert,liu2019roberta}, and domain pretraining helps financial sentiment \emph{classification} \citep{araci2019finbert}, yet their 512-token window falls more than an order of magnitude short of the median long-form filing. Proposed remedies---truncation, chunk pooling, hierarchical chunk encoders \citep{pappagari2019hierarchical}, and native sparse attention \citep{beltagy2020longformer}---are benchmarked almost exclusively on classification, leaving their behaviour on noisy financial \emph{regression} targets untested.

\paragraph{LLMs as financial forecasters.}
Prompted instruction LLMs have recently been proposed as market forecasters, eliciting predictions in natural language rather than fitting a head to learned features \citep{lopezlira2023chatgpt}. Two worries dominate: contamination---a model trained through part of the evaluation period may have memorised era- or entity-specific outcomes, so apparent skill reflects the training corpus rather than the input text---and elicitation fragility, since prompted forecasts can hinge on phrasing and decoding choices. We treat temperature-zero prompting of a 32B instruction model \citep{qwen3_2025} as a new point on the representation spectrum---the same filing text read by a generative reader rather than pooled into embeddings---and discipline it with date-only and date-plus-ticker contamination prompts, an era split at the training cutoff, paraphrase and repeat-decode checks, and replication tests in two further LLM families.

\paragraph{Benchmark hygiene and identity leakage.}
Beyond finance, benchmark gains are known to arise from shortcut features rather than the capability under test \citep{geirhos2020shortcut}, and entity memorisation is a canonical shortcut. Disclosure panels are acutely exposed: firms file repeatedly and volatility is firm-persistent, so an expressive representation can encode \emph{which firm is speaking} rather than \emph{what is said}, inflating apparent gains over any identity-blind reference. We claim no generality beyond this benchmark; our contribution is a protocol that makes the artefact measurable in this instance, by running a zero-text firm-identity forecast through the same combination machinery as every text model. To our knowledge, no prior work jointly imposes a survivorship-free point-in-time panel, a recalibrated reference interval---from a single HAR to a firm-identity-augmented reference, closed by a maximal price pool---day-clustered Diebold--Mariano inference with Holm correction and placebo gating, and a model spectrum from dictionaries to a prompted instruction LLM. Supplying that controlled test is the contribution of this paper.

